\documentclass{amsart}
% \numberwithin{equation}{section}  % removed: no sections, was producing "0.X" labels
\usepackage[utf8]{inputenc}
\usepackage[T1]{fontenc}
\usepackage{lmodern}
%\usepackage[mathscr]{euscript}
\usepackage[margin=1.2in]{geometry}
\usepackage{setspace}
\setstretch{1.125}
\usepackage[usenames, dvipsnames]{xcolor}
\usepackage{graphicx}
\usepackage{mathtools}
\usepackage{amssymb}
\usepackage{amsthm}
%\usepackage{xparse}
\usepackage{xspace}
\usepackage{enumitem}
\usepackage{tikz-cd}
\usepackage[backref=page, bookmarks=false]{hyperref}
\usepackage[capitalize, noabbrev]{cleveref}
\usepackage{dynkin-diagrams}
% \usepackage{tocloft}

% % Adjust the indentation for subsections in the table of contents
% \cftsetindents{subsection}{1.5em}{2.5em}


\setcounter{tocdepth}{2}

\newcommand{\Erdos}{Erdős\xspace}

\newcommand{\mb}{\mathbf}

\newcommand{\iso}{\cong}
\newcommand{\isoto}{\xrightarrow{\simeq}}

\newcommand{\Z}{\mathbb Z}
\newcommand{\Q}{\mathbb Q}
\newcommand{\R}{\mathbb R}
\newcommand{\C}{\mathbb C}
\newcommand{\T}{\mathbb T}
\newcommand{\F}{\mathbb{F}}
\newcommand{\ii}{\mathbf i}
\newcommand{\mcal}[1]{\mathcal{#1}}

\newcommand{\op}{\mathsf{op}}
\newcommand{\GL}{\mathrm{GL}}
\newcommand{\PSL}{\mathrm{PSL}}
% \newcommand{\Diff}{\mathrm{Diff}}
\newcommand{\Lie}{\mathrm{Lie}}
\newcommand{\Zn}{\Z(n)}
\newcommand{\Rn}{\R(n)}
\newcommand{\Bdot}{B_\bullet}
\newcommand{\Edot}{E_\bullet}
\newcommand{\BdotG}{B_\bullet G}
\newcommand{\EdotG}{E_{\bullet}G}
\newcommand{\BdotH}{B_\bullet H}
\newcommand{\BnablaG}{B_\nabla G}
\newcommand{\EnablaG}{E_\nabla G}
\newcommand{\BdotGamma}{B_\bullet \Gamma}
\newcommand{\GFG}{\Gamma \backslash F / G}
\newcommand{\GFR}{\Gamma \backslash F / \R}
\newcommand{\GF}{\Gamma \backslash F}
\newcommand{\Diff}{\mathrm{Diff}^+(S^1)}
\newcommand{\Mfld}{Mfld}
\newcommand{\Witt}{\mathfrak{w}}
\newcommand{\LieVir}{\mathfrak{vir}}
% \newcommand{\FGG}(F/G}

\newcommand{\point}{\ast}

\newcommand{\Zf}{\mb{Z}_2^F}
\newsavebox{\pullback}
\sbox\pullback{%
\begin{tikzpicture}%
\draw (0,0) -- (1ex,0ex);%
\draw (1ex,0ex) -- (1ex,1ex);%
\end{tikzpicture}}

\DeclareMathOperator{\Ext}{Ext}
\DeclareMathOperator{\Hom}{Hom}

\newtheorem{lemma}[equation]{Lemma}
\newtheorem{corollary}[equation]{Corollary}
\newtheorem{proposition}[equation]{Proposition}
\newtheorem{thm}[equation]{Theorem}
\newtheorem{theorem}[equation]{Theorem}
\newtheorem*{mainthm}{\cref{main_thm}}
\newtheorem{conjecture}[equation]{Conjecture}

\theoremstyle{definition}
\newtheorem{example}[equation]{Example}
\newtheorem{definition}[equation]{Definition}
\newtheorem{claim}[equation]{Claim}
\newtheorem{conj}[equation]{Conjecture}
\newtheorem{question}[equation]{Question}

\theoremstyle{remark}
\newtheorem{remark}[equation]{Remark}
\newtheorem*{fct}{Fact}
\newtheorem*{note}{Note}
\newtheorem{observation}[equation]{Observation}

\crefname{thm}{Theorem}{Theorems}
\crefname{lem}{Lemma}{Lemmas}
\crefname{cor}{Corollary}{Corollaries}
\crefname{prop}{Proposition}{Propositions}
\crefname{ex}{Exercise}{Exercises}
\crefname{exm}{Example}{Examples}
\crefname{defn}{Definition}{Definitions}
\crefname{claim}{Claim}{Claims}
\crefname{rem}{Remark}{Remarks}
\crefname{fct}{Fact}{Facts}
\crefname{note}{Note}{Notes}

\newcommand{\term}{\emph} % e.g. "The \term{trace} is defined to be..."

\DeclarePairedDelimiter\abs{\lvert}{\rvert}
\DeclarePairedDelimiter\set{\{}{\}}
\DeclarePairedDelimiter\paren{(}{)}
\DeclarePairedDelimiter\ang{\langle}{\rangle}
\DeclarePairedDelimiter\ket{\lvert}{\rangle}

\makeatletter
	\let\oldparen\paren
	\def\paren{\@ifstar{\oldparen}{\oldparen*}}
\makeatother

\usepackage{microtype}
\usepackage{hypcap}

\hypersetup{
 colorlinks,
 linkcolor={red!50!black},
 citecolor={green!50!black},
 urlcolor={blue!80!black}
}

% make cleveref use the Oxford comma:
% see https://tex.stackexchange.com/questions/161338
\newcommand{\creflastconjunction}{, and\nobreakspace}

\newcommand{\TODO}{\textcolor{red}{TODO}}
\newcommand{\QUESTION}{\textcolor{green}{QUESTION}}
\newcommand{\ADDFIGURE}{\textcolor{blue}{ADDFIGURE}}


\newcommand{\cat}{\mathsf}
\newcommand{\Man}{\cat{Man}}
\newcommand{\Sh}{\cat{Sh}}
\newcommand{\Ch}{\cat{Ch}}
\newcommand{\sSet}{\cat{sSet}}

\newcommand{\Hsm}{H_{\mathrm{SM}}}

\newcommand{\Sym}{\mathrm{Sym}}
\newcommand{\fg}{\mathfrak g}
\newcommand{\gl}{\mathfrak{gl}}
\newcommand{\tr}{\mathrm{tr}}
\newcommand{\su}{\mathfrak{su}}
\newcommand{\OmegaR}{\Omega^1 \otimes i\R}
\newcommand{\chern}{(c_1)^\nabla}


\renewcommand{\d}{\mathrm d}
\newcommand{\ud}{\,\d}

\newcommand{\tGam}{\widetilde\Gamma}
\newcommand{\CExt}{\mathrm{CExt}}
\newcommand{\Map}{\mathrm{Map}}

% was this deifned somewhere else...?
\newcommand{\AlgVir}{\mathfrak{vir}}
\newcommand{\GpVir}{\widetilde\Gamma}

% maintenance
\newcommand{\YL}[1]{\footnote{\color{orange} {\bf YL:} {#1}}}
\newcommand{\YLComm}[1]{{\color{orange} {\bf YL:} {#1}}}


\newcommand{\IZ}{I\mb{Z}}

\newcommand{\needcite}[1]{{\color{red} cite[{#1}]}}
\newcommand{\Question}{{\color{red} Question}}
\newcommand{\cF}{\mathcal{F}}
\newcommand{\cA}{\mathcal{A}}

\newcommand{\vecn}{\vec{n}}
\newcommand{\vecm}{\vec{m}}

\newcommand{\mbb}[1]{\mathbb{{#1}}}
\newcommand{\Moduli}{{\mbb{H}/SL(2, \mbb{Z})}}
%%%% CFT notations
% \newcommand{\tr}{\mathrm{tr}}


\newcommand{\LambdaR}{\Lambda_\R}
\newcommand{\LambdaRplus}{\LambdaR^+}
\newcommand{\LambdaRminus}{\LambdaR^-}
\newcommand{\U}{\mathrm{U}}
\newcommand{\cO}{\mathcal{O}}
\newcommand{\NLambda}{N_{\Lambda}}
\newcommand{\vol}{\mathrm{vol}}
\newcommand{\QL}{Q^L}
\newcommand{\QR}{Q^R}
% \renewcommand{\H}{\mathbb{H}}  % removed: conflicts with Hungarian umlaut needed for "Erdős"
\newcommand{\ETS}{E^{\mathrm{TS}}}
\newcommand{\odd}{\mathrm{odd}}
\newcommand{\even}{\mathrm{even}}
\newcommand{\cusp}{\mathrm{cusp}}
\newcommand{\TS}{\mathrm{TS}}
\newcommand{\SL}{\mathrm{SL}}
\newcommand{\Mp}{\mathrm{Mp}}
\renewcommand{\Re}{\mathrm{Re}}
\renewcommand{\Im}{\mathrm{Im}}
\newcommand{\tZ}{\tilde{Z}}
\newcommand{\Ann}{\mathrm{Ann}}
\newcommand{\Res}{\mathrm{Res}}
\newcommand{\Weil}{\mathrm{Weil}}
\newcommand{\cE}{\mathcal{E}}
% \newcommand{\PSL}{\mathrm{PSL}}
\newcommand{\res}{\mathrm{res}}
\newcommand{\SO}{\mathrm{SO}}
\newcommand{\triv}{\mathrm{triv}}
\newcommand{\Zctimes}[1]{(\Z/{#1}\Z)^\times}
\newcommand{\Zmodc}[1]{\Z{#1}\Z}
\newcommand{\lcm}{\mathrm{lcm}}
\newcommand{\CP}{\mathbb{C}P}
\newcommand{\cond}{\mathrm{cond}}
\newcommand{\cC}{\mathcal{C}}
\newcommand{\cD}{\mathcal{D}}
\newcommand{\Rep}{\mathrm{Rep}}
\newcommand{\Mpn}[1]{\Mp^{({#1})}_2(\Z)}
\newcommand{\Sp}{\mathrm{Sp}}
\newcommand{\diag}{\mathrm{diag}}
\renewcommand{\sp}{\mathfrak{sp}}
\newcommand{\mfh}{\mathfrak{h}}
\newcommand{\mfg}{\mathfrak{g}}
\newcommand{\lrangle}[1]{\langle #1 \rangle}

% Sep-3-25

\newcommand{\N}{\mathbb{N}}
\newcommand{\Aut}{\mathrm{Aut}}
\renewcommand{\P}{\mathbb{P}}
\newcommand{\OF}{\mathcal{O}_F}
\renewcommand{\OE}{\mathcal{O}_E}
\newcommand{\mD}{\mathfrak{D}}
\newcommand{\mq}{\mathfrak{q}}
\newcommand{\ma}{\mathfrak{a}}
\newcommand{\OK}{\mathcal{O}_K}
\newcommand{\cM}{\mathcal{M}}
\newcommand{\SU}{\mathrm{SU}}
\newcommand{\real}{\mathrm{real}}
\newcommand{\complex}{\mathrm{complex}}
\newcommand{\Cl}{\mathrm{Cl}}
\newcommand{\PGL}{\mathrm{PGL}}
\newcommand{\Spec}{\mathrm{Spec}}

% ---------- Notation ----------
\newcommand{\A}{\mathbb{A}}
\newcommand{\Asai}{\mathrm{Asai}}
\newcommand{\Ad}{\mathrm{Ad}}
\newcommand{\BC}{\mathrm{BC}}   % (stable) base change
\newcommand{\AI}{\mathrm{AI}}   % automorphic induction
\newcommand{\zetaE}{\zeta_E}
\newcommand{\zetaF}{\zeta_F}
\newcommand{\LambdaE}{\Lambda_E}
\newcommand{\LambdaF}{\Lambda_F}
\newcommand{\cS}{\mathcal{S}}
\newcommand{\cR}{\mathcal{R}}   % Rankin–Selberg transform
\newcommand{\Tr}{\mathrm{Tr}}
\newcommand{\Gal}{\mathrm{Gal}}
\renewcommand{\abs}[1]{\left\lvert #1 \right\rvert}
\newcommand{\norm}[1]{\left\lVert #1 \right\rVert}
\newcommand{\veps}{\varepsilon}
% ---------- Groups/Weil ----------
\newcommand{\WE}{W_E}
\newcommand{\WF}{W_F}
\newcommand{\WEd}{W_E'}
\newcommand{\WFd}{W_F'}   % Weil–Deligne groups
% ---------- Tensor induction symbol ----------
\newcommand{\Indten}{\mathrm{Ind}^{\,\otimes}} % tensor (a.k.a. multiplicative) induction
\newcommand{\Rplus}{\R_{>0}}
\newcommand{\nc}{\mathrm{nc}}

\title[\Erdos Problem \#675(a)]{\Erdos Problem \#675, subquestion (a): \\ {\tiny a super-polynomial lower bound on local periods of the sums-of-two-squares set}}
\author{Yu Leon Liu}
\address{1 Oxford St, Cambridge, MA 02139}
\email{yuleonliu@math.harvard.edu}

\begin{document}

\begin{abstract}
Let $E$ be the set of sums of two squares. We prove that, for every fixed
$0 < c < 1/10$ and all sufficiently large $n$, every positive integer $t$
satisfying the $n$-translation property for $E$ (that is, $a \in E
\Longleftrightarrow a + t \in E$ for all $1 \leq a \leq n$) must exceed
$\exp(n^c)$. We do not address the existence of such $t$, that is,
whether $E$ has the translation property.
\end{abstract}

    \maketitle

\begin{definition}\label{def:trans}
    For a subset $A \subset \N$ and positive integer $n$, we say that a positive integer $t$ satisfies the $n$-translation property for $A$ if
    \[
    a\in A \quad \Longleftrightarrow \quad a+t \in A
    \]
    holds for all $1 \leq a \leq n$. We say that $A$ has the \textit{translation property} if for every $n > 0$ there exists $t$ satisfying the $n$-translation property for $A$.
\end{definition}

Problem 675~\cite{bloom675} has three parts:
\begin{quote}
\begin{itemize}[leftmargin=1.2em]
\item[(a)] Does the set of sums of two squares have the translation property?
\item[(b)] If we partition all primes into $P\sqcup Q$ such that each part contains $\gg x/\log x$ many primes $\le x$ for all large $x$, does the set of integers only divisible by primes from $P$ have the translation property?
\item[(c)] If $A$ is the set of squarefree numbers, how fast does the minimal such $t_n$ satisfying the $n$-translation property of $A$ grow? Is $t_n>\exp(n^c)$ for some constant $c>0$?
\end{itemize}
\end{quote}

Here we make the following progress on part (a) concerning sums of two squares:
\begin{theorem}\label{thm:main-thm}
    Let $E$ be the set of positive integers expressible as $a^2 + b^2$ with $a, b \in \mathbb Z_{\geq 0}$. For every $0 < c < \frac{1}{10}$, there exists $N_c$ such that for all $n \geq N_c$, every $t$ satisfying the $n$-translation property for $E$ satisfies
\[
    t > \exp(n^c).
\]
\end{theorem}
\begin{remark}
    We \emph{do not} prove part (a), i.e., that the set $E$ satisfies the translation property. We do show that if $E$ has the translation property, then $t$ grows faster than $\exp(n^c)$ as in \cref{thm:main-thm}.
\end{remark}
\begin{remark}
    We discuss potential improvements on the constant $\frac{1}{10}$ in \cref{rem:end}.
\end{remark}

First we note that 
\begin{equation}\label{eq:E-p-3}
    E = \{m \geq 1 \mid v_p(m) \text{ is even for all primes } p \equiv 3 \pmod 4\}.
\end{equation}
We have the following lemma:
\begin{lemma}\label{lem:1}
    Let $p \equiv 3 \pmod 4$ and $n \geq 1$. Suppose $t$ satisfies the $n$-translation property for $E$. Then for any $a\in E \cap [1,n]$ with
    \begin{equation*}
        a \equiv  -t \pmod p,
    \end{equation*}
    we have
    \begin{equation*}
        a \equiv -t \pmod{p^2}.
    \end{equation*}
\end{lemma}
\begin{proof}
    By construction we have $p \mid a + t$, and the $n$-translation property implies that $a + t \in E$. It follows from \eqref{eq:E-p-3} that $p^2 \mid a + t$, and thus
    \begin{equation*}
        a  \equiv -t \pmod{p^2}.
    \end{equation*}
\end{proof}

Next, we show that any sufficiently small prime $p \equiv 3 \pmod 4$ must divide $t$. The idea is to find $a,b$ that are both $\equiv -t \pmod p$ but differ modulo $p^2$.
\begin{lemma}\label{lem:2}
   Let $p \equiv 3 \pmod 4$ be prime. There exists $N_p$ such that,
   for every $n \geq N_p$ and every $t \geq 1$ satisfying the $n$-translation property for $E$, we have $p \mid t$.
\end{lemma}

\begin{proof}
    For each nonzero residue class $r \pmod p$, the residues $r$ and $r + p$ differ modulo $p^2$. By the Chinese remainder theorem, there are unique classes $u_{r,0}, u_{r,1} \pmod{4p^2}$ satisfying
\begin{align*}
u_{r,0}&\equiv 1\pmod 4,\quad u_{r,0}\equiv r\pmod{p^2},\\
u_{r,1}&\equiv 1\pmod 4,\quad u_{r,1}\equiv r+p\pmod{p^2}.
\end{align*}
    Since $r \not\equiv 0 \pmod p$ and each $u_{r,j} \equiv 1 \pmod 4$, both $u_{r,0}$ and $u_{r,1}$ are coprime to $4p^2$.
    By Dirichlet's theorem on primes in arithmetic progressions,
    there exist primes $\ell_{r,0}$ and $\ell_{r,1}$ with
\[
\ell_{r,0}\equiv u_{r,0}\pmod{4p^2}, \qquad \ell_{r,1}\equiv u_{r,1}\pmod{4p^2}.
\]
Each $\ell_{r,j} \equiv 1 \pmod 4$, hence lies in $E$ by \eqref{eq:E-p-3}. Let
\begin{equation*}
    N_p \coloneqq \max_{\substack{1 \leq r \leq p-1 \\ j \in \{0,1\}}} \ell_{r,j}.
\end{equation*}

Now suppose $n \geq N_p$ and $t$ has the $n$-translation property for $E$, and assume for contradiction that $p \nmid t$. Pick $r \in \{1, \dots, p-1\}$ with $r \equiv -t \pmod p$. The witnesses $\ell_{r,0}, \ell_{r,1} \in E \cap [1,n]$ both satisfy $\ell_{r,j} \equiv r \equiv -t \pmod p$, so \cref{lem:1} forces $\ell_{r,j} \equiv -t \pmod{p^2}$ for $j \in \{0,1\}$. In particular $\ell_{r,0} \equiv \ell_{r,1} \pmod{p^2}$. But by construction $\ell_{r,0} \equiv r \pmod{p^2}$ and $\ell_{r,1} \equiv r + p \pmod{p^2}$, so $\ell_{r,0} \not\equiv \ell_{r,1} \pmod{p^2}$, a contradiction. It follows that $p \mid t$.
\end{proof}

\begin{remark}
    In the proof above we chose $\ell_{r,j}$ to be primes because that guarantees they lie in $E$. In fact, any choice of $\ell_{r,0}, \ell_{r,1} \in E$ satisfying
    \begin{equation*}
        \ell_{r,0} \equiv r \pmod{p^2}, \qquad \ell_{r,1} \equiv r + p \pmod{p^2}
    \end{equation*}
    suffices, and a careful choice could yield a better bound.
\end{remark}

Next, we bound $N_p$ as a function of $p$:
\begin{theorem}[Linnik~\cite{linnik1944}]\label{thm:A}
   There exist constants $C > 0$ and $L > 0$ such that for any modulus $q \geq 2$ and any residue class $a$ with $1 \leq a \leq q-1$ and $\gcd(a,q) = 1$, the smallest prime $p_0 \equiv a \pmod q$ satisfies $p_0 \leq C q^L$.
\end{theorem}
The best current unconditional value is $L \leq 5$, due to Xylouris~\cite{xylouris2011}.

\begin{proof}[Proof of \cref{thm:main-thm}]
    Let $p \equiv 3 \pmod 4$ be a prime.
    Applying \cref{thm:A} with $q = 4p^2$, we see that we can choose primes $\ell_{r,j}$ with
    \[ \ell_{r,j} \leq C (4p^2)^L.\]
    Thus, with $A = 4^L C$, we have
    \begin{equation*}
        N_p \leq A p^{2L}.
    \end{equation*}
    In particular, if $p \leq (n/A)^{1/(2L)}$, then $n \geq N_p$, so any $t$ satisfying the $n$-translation property for $E$ has $p \mid t$ by \cref{lem:2}.

    Let $y \coloneqq (n/A)^{1/(2L)}$. It follows that
    \[
\prod_{\substack{p\le y\\ p\equiv 3 \,(\bmod\, 4)}} p \;\Big|\; t.
\]
By the prime number theorem in arithmetic progressions (see e.g.\ Davenport~\cite{davenport2000}), the primes $\equiv 3 \pmod 4$ have density $\tfrac12$ among the odd primes, so
\begin{equation}\label{eq:pnt-ap}
    \sum_{\substack{p\le y\\ p\equiv 3 \,(\bmod\, 4)}} \log p \;\sim\;
    \frac{y}{2} \qquad \text{as } y \to \infty.
\end{equation}
Fix a constant $c_1 \in (0, \tfrac12)$, say $c_1 = \tfrac14$. By \eqref{eq:pnt-ap}, there exists a threshold $y_0 = y_0(c_1)$ such that for all $y \geq y_0$,
\begin{equation}\label{eq:pnt-bound}
    \sum_{\substack{p\le y\\ p\equiv 3 \,(\bmod\, 4)}} \log p \;\geq\; c_1 y.
\end{equation}
Set $N_1 \coloneqq A \cdot y_0^{2L}$, so that $n \geq N_1$ guarantees $y = (n/A)^{1/(2L)} \geq y_0$. For all such $n$, combining the divisibility above with \eqref{eq:pnt-bound} gives
\begin{equation*}
    \log t \;\geq\; \sum_{\substack{p\le y\\ p\equiv 3 \,(\bmod\, 4)}} \log p \;\geq\; c_1 y \;=\; K n^{1/(2L)},
\end{equation*}
where $K \coloneqq c_1 A^{-1/(2L)}$.

Now fix any $0 < c < 1/(2L)$; since $L \leq 5$, every $0 < c < 1/10$ satisfies this. Since the exponent gap $1/(2L) - c$ is positive, the inequality $K n^{1/(2L)} > n^c$, equivalently $n^{1/(2L) - c} > K^{-1}$, holds for all $n \geq N_2(c) \coloneqq \lfloor K^{-2L/(1 - 2Lc)} \rfloor + 1$. Setting
\[
N_c \coloneqq \max\bigl(\lceil N_1 \rceil,\; N_2(c)\bigr),
\]
we conclude that for every $n \geq N_c$,
\[
\log t \;\geq\; K n^{1/(2L)} \;>\; n^c, \qquad \text{i.e., } t > \exp(n^c). \qedhere
\]
\end{proof}

\begin{remark}\label{rem:end}
The constant $\tfrac{1}{10}$ in \cref{thm:main-thm} arose as $\tfrac{1}{2L}$ with $L \leq 5$ from \cite{xylouris2011}. There are two potential ways to improve it.

\smallskip
\noindent\emph{First, by improving Linnik's constant.} Under the Generalized Riemann Hypothesis, one has $p_{\min}(a, q) \ll \varphi(q)^2 (\log q)^2$, which sets Linnik's constant to $L \leq 2 + \epsilon$ in \cref{thm:A}. This pushes the conclusion of \cref{thm:main-thm} to $t > \exp(n^c)$ with $0 < c < \tfrac{1}{4}$. The Linnik conjecture, $p_{\min}(a, q) \ll_\varepsilon q^{1+\varepsilon}$ for every $\varepsilon > 0$, would further yield $t > \exp(n^c)$ for every $0 < c < \tfrac{1}{2}$.

\smallskip
\noindent\emph{Second, by relaxing the choice of witnesses.} As noted in the remark following \cref{lem:2}, we never used that the witnesses $\ell_{r,j}$ are prime; any $\ell_{r,j} \in E$ with the prescribed residue $\pmod{p^2}$ suffices. Bounding the smallest such $\ell_{r,j}$ is a question about $E$ in arithmetic progressions; sharp bounds here would likely improve the constant further.
\end{remark}

\section*{Acknowledgements}
The proof in this note was discovered with the assistance of OpenAI's
Codex together with the Rethlas open-source agentic
mathematics-research pipeline. The strategy is directly inspired by
Ho Boon Suan's proof of the analogous super-polynomial lower bound for
the squarefree subquestion (c) of \Erdos Problem 675~\cite{hobs2026}.
All mathematical arguments and claims in the final manuscript were
independently verified by the author, who takes full responsibility
for the paper.

\begin{thebibliography}{9}

\bibitem{bloom675}
T.~F.~Bloom,
\emph{\Erdos Problem \#675},
\url{https://www.erdosproblems.com/675},
accessed 10 May 2026.

\bibitem{hobs2026}
Ho Boon Suan,
\emph{A squarefree lower bound for \Erdos Problem 675},
note, 18 April 2026,
\url{https://boonsuan.github.io/erdos675_squarefree.pdf}.

\bibitem{davenport2000}
H.~Davenport,
\emph{Multiplicative Number Theory}, 3rd ed.,
Graduate Texts in Mathematics, vol.~74, Springer-Verlag, New York, 2000.
Revised and with a preface by Hugh L.~Montgomery.

\bibitem{linnik1944}
U.~V.~Linnik,
\emph{On the least prime in an arithmetic progression. I. The basic theorem},
Rec. Math. [Mat. Sbornik] N.S. \textbf{15(57)} (1944), no.~2, 139--178.

\bibitem{xylouris2011}
T.~Xylouris,
\emph{\"{U}ber die Nullstellen der Dirichletschen $L$-Funktionen und die kleinste Primzahl in einer arithmetischen Progression},
PhD thesis, Rheinische Friedrich-Wilhelms-Universit\"{a}t Bonn, 2011;
also published as \emph{Bonner Mathematische Schriften} \textbf{404}, Universit\"{a}t Bonn, Mathematisches Institut, Bonn, 2011.

\end{thebibliography}

\end{document}
